\SourceLeaf{207}\clearpage
\begin{center}
[ÜBER DIE ERSTEN GRÜNDE DER GEOMETRIE.]
\end{center}

\begin{center}
\textsc{Gauss} an \textsc{Bessel}. Göttingen, 27. Januar 1829.
\end{center}

.\ .\ .\ .\ .\ .\ .\ Auch über ein anderes Thema, das bei mir schon fast 40 Jahr
alt ist, habe ich zuweilen in einzelnen freien Stunden wieder nachgedacht, ich
meine die ersten Gründe der Geometrie: ich weiss nicht, ob ich Ihnen je über
meine Ansichten darüber gesprochen habe. Auch hier habe ich manches noch
weiter consolidirt, und meine Überzeugung, dass wir die Geometrie nicht vollständig a priori begründen können, ist, wo möglich, noch fester geworden.
Inzwischen werde ich wohl noch lange nicht dazu kommen, meine \textls{sehr ausgedehnten} Untersuchungen darüber zur öffentlichen Bekanntmachung auszuarbeiten, und vielleicht wird diess auch bei meinen Lebzeiten nie geschehen,
da ich das Geschrei der Böotier scheue, wenn ich meine Ansicht ganz aussprechen wollte. --- Seltsam ist es aber, dass \textls{ausser} der bekannten Lücke in
Euklids Geometrie, die man bisher umsonst auszufüllen gesucht hat, und nie
ausfüllen wird, es noch einen andern Mangel in derselben gibt, den meines
Wissens niemand bisher gerügt hat, und dem abzuhelfen keinesweges leicht
(obwohl möglich) ist. Diess ist die Definition des \textls{Planum} als einer Fläche,
in der die, \textls{irgend zwei} Punkte verbindende, gerade Linie \textls{ganz} liegt. Diese
Definition enthält \textls{mehr}, als zur Bestimmung der Fläche nöthig ist, und involvirt \emph{tacite} ein \textls{Theorem}, welches erst bewiesen werden muss. .\ .\ .\ .\ .\ .

\rule{4cm}{0.5pt}

